\documentclass[10pt,a4paper]{amsart}
\usepackage[margin=19mm]{geometry}
\usepackage{amsmath,amssymb,mathtools}
\usepackage[scr=rsfs,cal=euler]{mathalfa}
\usepackage{xfrac}
\usepackage[unicode,hidelinks,bookmarks=false]{hyperref}
\newcommand{\ie}{i.e.\ }
\newcommand{\cf}{cf.\ }
\newcommand{\resp}{respectively\ }
\newcommand{\Subsection}{\subsection}
\providecommand{\nobreakdash}{\nobreak\hbox{-}}
\setlength{\emergencystretch}{2em}
\multlinegap0pt
\usepackage{bm}
\usepackage{bbm}
\usepackage{dsfont}
\usepackage{xspace}
\usepackage{cancel}
\usepackage{mathtools}
\usepackage{xcolor}
\usepackage{amsthm}
\makeatletter
\@ifundefined{Theorem}{\newtheorem{Theorem}{Theorem}}{}
\@ifundefined{Lemma}{\newtheorem{Lemma}[Theorem]{Lemma}}{}
\makeatother
\def\R{\mathbb{R}}
\newcommand{\De} {\Delta}
\title[Symmetry and Approximate Symmetry for Solutions of Mixed Loc]{Symmetry and Approximate Symmetry for Solutions of Mixed Local-Nonlocal Singular Equations}
\author{Sanjit Biswas}
\date{Source: arXiv:2602.19054v2 (CC BY 4.0)}
\begin{document}
\maketitle

\noindent\emph{Source and licence.} Symmetry and Approximate Symmetry for Solutions of Mixed Local-Nonlocal Singular Equations. Version arXiv:2602.19054v2 (CC BY 4.0), distributed under the Creative Commons Attribution 4.0 International licence, as recorded in the arXiv licence metadata for this exact version and shown on the abstract page https://arxiv.org/abs/2602.19054v2. Full rights notice: licenses/arxiv-2602.19054-CC-BY-4.0.txt.\par\smallskip

\noindent\textit{Editorial scope.} Counted unit: the Lemma whose source label is \texttt{SML} in the article (source0.tex lines 540--550, source label \texttt{SML}), delivered together with the article's own complete original proof of it (source0.tex lines 551--613, one \texttt{proof} environment of 63 lines). No mathematics is written, repaired or re-derived: statement and proof are the article's own text line for line. Required context retained uncounted: Sa2 (lines 529--535). Every definition, display and label of those lines is retained unchanged. Omitted from this package and not redistributed: 1--528, 536--539, 614--1041. Nothing inside a retained excerpt is omitted or altered: every retained body is delivered exactly as the article's lines are written, so no omission receipt of any kind accompanies this package. Citations: the retained excerpts carry no citation command at all, so no bibliography entry of the article is transcribed and no reference is redistributed. Reference adaptation: none was needed and none was made; every reference inside the delivered unit resolves inside it, so no unresolved reference or citation is produced. Printed slips kept literal: no printed word, symbol, hypothesis or number of the counted statement or of its proof was altered. Layout and class: the article's own class is \texttt{article} with packages including amssymb, amsmath, comment, amsthm, colortbl, hyperref, eucal, epsf, esint, a4wide, xy; the wrapper is the lane's pinned \texttt{amsart} editorial wrapper at 10pt on A4, which restates the theorem-like environments the retained text needs and adds the article's own macro definitions that the retained text needs (\texttt{\textbackslash De}, \texttt{\textbackslash R}), with extra wrapper packages (bm, bbm, dsfont, xspace, cancel, mathtools, xcolor). The delivered numbering is the unit's own. Assets: no retained span contains a figure environment, a graphics-inclusion command, a PDF-page inclusion, a table or a TikZ picture; no image file is copied into this package, the package declares no asset dependency and no image file is redistributed with it. Licence: version arXiv:2602.19054v2 of this article is distributed under the Creative Commons Attribution 4.0 International licence, as recorded in the arXiv licence metadata for this exact version and shown on the abstract page https://arxiv.org/abs/2602.19054v2. A full rights notice is delivered at licenses/arxiv-2602.19054-CC-BY-4.0.txt. Characters: every retained excerpt is pure ASCII, so the delivered engine, XeLaTeX with its default Latin Modern text font, reports no missing character.
\begin{Lemma}\label{Sa2}
    Suppose that $p,\,q\geq 1$ are two real numbers such that $2(q-1)=(p-1)$. Then for all $t$, the following holds:
    \begin{itemize}
        \item[(a)] $\xi'(t)=\eta'(t)^2$,
        \item[(b)] $t\xi(t)\leq q\eta^2(t)$.
    \end{itemize}     
\end{Lemma}

\begin{Lemma}\label{SML}
     Let $\Omega\subset\R^n$ be a bounded Lipschitz domain. Suppose that $u\in C^1 (\overline{\Omega})\textcolor{black}{\cap \mathcal{L}^s(\R^n)}$ is a weak subsolution of the equation
\begin{align}\label{Sa}
    \begin{cases}
        -\De u+(-\De)^su+c(x)u=g \text{ in }\Omega,\\
        u\leq 0\text{ on }\textcolor{black}{\R^n\setminus\Omega},
    \end{cases}
\end{align}
where $c,\, g\in L^\infty(\Omega).$ Then, there exist two constants $\delta=\delta(n,\|c\|_{L^\infty(\Omega)})>0$ and $\mathcal{K}=\mathcal{K}(n,\|c\|_{L^\infty(\Omega)})>0$
such that whenever $|\Omega|\leq\delta$, we have $$\|u^+\|_{L^\infty(\Omega)}\leq \mathcal{K} \|g\|_{L^\infty(\Omega)}.$$    
\end{Lemma}

\begin{proof}
Without loss of generality, we assume that $|\Omega|\leq 1$ and $n>2$. We define $w:=u^++k$, where $k=\|g\|_{L^\infty(\Omega)}$. We assume $l>k$ in the aforementioned definition of $\xi$ and $\eta$. Incorporating $\phi=\xi(w)$ in the weak formulation of \eqref{Sa}, we obtain 
\begin{align*}
    \int_\Omega \xi'(w)\nabla u\cdot\nabla w\;dx+&\int_{\R^n}\int_{\R^n}\frac{(u(x)-u(y))(\xi(w)(x)-\xi(w)(y))}{|x-y|^{n+2s}}\;dy\;dx\nonumber\\
    &+\int_{\Omega}c(x)u(x)\xi(w)(x)\;dx\leq\int_{\Omega}g\xi(w)\;dx.
\end{align*}
Using the monotonicity property of $\xi$ and the fact that $ u\geq 0$ in $\mathrm{supp}(\xi(w))$, we obtain
\begin{align*}
    \int_\Omega \xi'(w)|\nabla w|^2\;dx+\int_{\Omega}c(x)u^+\xi(w)\;dx\leq\int_{\Omega}g\xi(w)\;dx.
\end{align*}
Utilizing Lemma \ref{Sa2}, one has
\begin{align}\label{Sa3}
     \int_\Omega |\nabla\eta(w)|^2\;dx\leq\int_{\Omega}(g-cu^+)\xi(w)\;dx&\leq (1+\|c\|_{L^\infty(\Omega)})\int_{\Omega}(k+u^+)\xi(w)\;dx \nonumber\\
     &\leq (1+\|c\|_{L^\infty(\Omega)})\int_\Omega w\xi(w)\;dx\nonumber\\
     &\leq q(1+\|c\|_{L^\infty(\Omega)})\int_\Omega |\eta(w)|^2\;dx.
\end{align}
Denote $v:=\eta(w)$. Since $q>1$, \eqref{Sa3} implies
\begin{align}\label{Sa4}
    \int_\Omega |\nabla v|^2\;dx\leq (1+\|c\|_{L^\infty(\Omega)})q^2\int_\Omega |v|^2\;dx.
\end{align}
Now, Using (\ref{Sa4}), Sobolev inequality and the fact $v>l^q$, we have
\begin{align*}
    \|v\|_{L^{2^*(\Omega)}}\leq \|v-l^q\|_{L^{2^*(\Omega)}}+\|l^q\|_{L^{2^*(\Omega)}}&\leq C\|\nabla v\|_{L^2(\Omega)}+\|l^q\|_{L^{2(\Omega)}}|\Omega|^{-\frac{1}{n}}\nonumber\\
    &\leq |\Omega|^{-\frac{1}{n}}\Big(Cq\|v\|_{L^2(\Omega)}+\|l^q\|_{L^2(\Omega)}\Big)\nonumber\\
    &\leq Cq|\Omega|^{-\frac{1}{n}}\|v\|_{L^2(\Omega)},
\end{align*}
where $C=C(n,\|c\|_{L^\infty(\Omega)})>1$ and $2^*:=\frac{2n}{n-2}$. Since $v\leq w^q+l^q-k^q$, we get
\begin{align*}
    \Big(\int_{l\leq w\leq r}v^{2^*}\;dx\Big)^\frac{1}{2^*}\leq Cq|\Omega|^{-\frac{1}{n}}\left(\|w^q\|_{L^2(\Omega)}+(l^q-k^q)|\Omega|^\frac{1}{2}\right).
\end{align*}
Taking $l\to k$ and $r\to\infty$, we deduce
\begin{align*}
    \|w\|_{L^{2^*q}(\Omega)}^q\leq Cq|\Omega|^{-\frac{1}{n}}\|w\|_{L^{2q}(\Omega)}^q.
\end{align*}
This inequality yields that for $\chi=\frac{n}{n-2}$, one has
\begin{align*}
    \|w\|_{L^{2q\chi}(\Omega)}\leq (Cq)^\frac{1}{q}|\Omega|^{-\frac{1}{nq}}\|w\|_{L^{2q}(\Omega)}.
\end{align*}
 Iterating this process, for every $i\in\mathbb{N}$ we obtain
\begin{align*}
    \|w\|_{L^{2\chi^i}(\Omega)}\leq (C|\Omega|^{-\frac{1}{n}})^{1+\frac{1}{\chi}+..+\frac{1}{\chi^{i-1}}}(\prod_{j=1}^{i-1} \chi^\frac{j}{\chi^j})\|w\|_{L^{2}(\Omega)}\leq C|\Omega|^{-\frac{1}{2}}\|w\|_{L^{2}(\Omega)},
\end{align*}
for some constant $C=C(n,\|c\|_{L^\infty(\Omega)})>0$. Taking $i\to\infty$, we deduce 
\begin{align*}
    \|w\|_{L^\infty(\Omega)}\leq C|\Omega|^{-\frac{1}{2}}\|w\|_{L^{2}(\Omega)}.
\end{align*}
Since $w=u^++k$, the above inequality reveals 
\begin{align}\label{Sa5}
    \|w\|_{L^\infty(\Omega)}\leq C|\Omega|^{-\frac{1}{2}}\|u^+\|_{L^2(\Omega)}+Ck,
\end{align}
for some constant $C=C(n,\|c\|_{L^\infty(\Omega)})>0$.
Moreover, putting $u^+\in H^1_0(\Omega)$ in the weak formulation of \eqref{Sa} and utilizing the Poincar\'{e} inequality, one has
\begin{align}\label{Sa6}
    \|\nabla u^+\|_{L^2(\Omega)}^2\leq \int_\Omega (g-cu)u^+\;dx\leq (1+\|c\|_{L^\infty(\Omega)})\int_\Omega wu^+\;dx\leq C\|w\|_{L^\infty(\Omega)}^2,
\end{align}
Combining \eqref{Sa5}, (\ref{Sa6}) together with the poincar\'{e} inequality, we get
\begin{align}\label{Sa7}
    \|w\|_{L^\infty(\Omega)}\leq C|\Omega|^{\frac{n-1}{2}}\|w\|_{L^\infty(\Omega)}+Ck,
\end{align}
where $C=C(n,\|c\|_{L^\infty(\Omega)})>0$ is a positive constant. Whenever $C|\Omega|^\frac{n-1}{2}\leq \frac{1}{2}$ i.e., $|\Omega|\leq \Big(\frac{1}{2C}\Big)^\frac{2}{n-1}(:=\delta)$, we have $$\|u^+\|_{L^\infty(\Omega)}\leq \|w\|_{L^\infty(\Omega)}\leq 2Ck=2C\|g\|_{L^\infty(\Omega)},$$
for some positive constant $C=C(n,\|c\|_{L^\infty(\Omega)}).$
This completes the proof.
\end{proof}

\end{document}
